% Generated draft: requires mathematical and source review.
\documentclass[11pt]{article}
\usepackage[margin=25mm]{geometry}
\usepackage{fontspec}
\setmainfont{Latin Modern Roman}
\usepackage{amsmath,amssymb,mathtools,unicode-math}
\setmathfont{Latin Modern Math}
\usepackage{enumitem,needspace,xurl,hyperref}
\setlength{\parindent}{0pt}
\setlength{\parskip}{5pt}
\setlength{\emergencystretch}{3em}
\newcommand{\N}{\mathbb{N}}
\newcommand{\Z}{\mathbb{Z}}
\newcommand{\Q}{\mathbb{Q}}
\newcommand{\R}{\mathbb{R}}
\newcommand{\C}{\mathbb{C}}
\newcommand{\F}{\mathbb{F}}
\newcommand{\A}{\mathbb{A}}
\newcommand{\PP}{\mathbb{P}}
\newcommand{\E}{\mathbb{E}}
\newcommand{\eps}{\varepsilon}
\newcommand{\dd}{\,\mathrm{d}}
\begin{document}
\section*{UnsolvedMath ID: 30002865}
\textbf{Source problem code:} catalogue entry
\subsection{Definitions and mathematical statement}
Let \(d\geq 1\) be an integer and let \(R=\C[x_1,\ldots,x_d]\). An ideal projector is a \(\C\)-linear map \(P:R\to R\) satisfying \(P^2=P\) whose kernel \(I=\ker P\) is an ideal. Here zero-dimensionality refers to the quotient ring: \(\dim(R/I)=0\), not to the vector-space dimension of \(I\). For a proper ideal in this polynomial ring, this implies that \(n=\dim_{\C}(R/I)\) is finite and positive. Set \(V=\operatorname{im}P\); then \(R=I\oplus V\) and \(\dim_{\C}V=n\).

For precision, use the following fixed-range interpretation of Hermite approximation. Given \(n\) distinct points \(a_1,\ldots,a_n\in\C^d\), suppose the evaluation map
\[
V\longrightarrow\C^n,\qquad v\longmapsto\bigl(v(a_1),\ldots,v(a_n)\bigr)
\]
is an isomorphism. The corresponding Lagrange projector \(L:R\to V\) is the unique linear map satisfying
\[
(Lp)(a_j)=p(a_j)\qquad(p\in R,\ 1\leq j\leq n).
\]
Call \(P\) a Hermite projector if there is a sequence of such Lagrange projectors \(L_k:R\to V\), each using \(n\) distinct points, such that
\[
\forall p\in R,\qquad L_kp\longrightarrow Pp\quad\text{in the finite-dimensional vector space }V.
\]
Convergence in \(V\) means convergence of coordinates in any basis of \(V\). The supplied record does not reproduce the source's definitions of Hermite projector or its convergence convention; this fixed-range interpretation requires confirmation against the cited passage.

The question is whether, for every \(d\geq1\) and every ideal projector \(P\) as above,
\[
\left[\dim(R/\ker P)=0\ \text{ and }\ \ker P=(f_1,\ldots,f_d)\ \text{for some }f_1,\ldots,f_d\in R\right]
\quad\Longrightarrow\quad
P\text{ is a Hermite projector}.
\]
The original formulation calls the kernel a complete intersection; the supplied corrected formulation specifies a zero-dimensional ideal generated by \(d\) polynomials. No reducedness or simplicity of the common zeros is assumed.

The supplied dated curation labels the problem partially solved and states that the implication has been proved under an additional no-roots-at-infinity hypothesis, while the unrestricted implication remains unproved. That additional hypothesis is not specified precisely in the record and is not imposed on the question above. These are supplied literature assessments, not independently verified mathematical or current-status claims.
\subsection{Short English statement}
Must every ideal projector on \(\C[x_1,\ldots,x_d]\) whose kernel is a zero-dimensional complete intersection be approximable by polynomial interpolation projectors at distinct points? The draft makes the approximation convention explicit, but that convention needs source confirmation. The supplied curation reports a result with an additional no-roots-at-infinity hypothesis and leaves the unrestricted question unresolved.
\subsection{Sources}
\begin{itemize}
\item Local UnsolvedMath registry: supplied complete record for permanent ID \(30002865\), title ``Complete-Intersection Kernels of Ideal Projectors,'' source problem code \texttt{catalogue entry}, in the Oberwolfach Reports Open Problems collection. The registry records proposed year \(2015\) and supplies no named proposer.
\item Supplied source citation: \emph{Multivariate Splines and Algebraic Geometry (2016)}, Oberwolfach Reports, pp.\ 1139--1200, DOI \texttt{10.4171/owr/2015/21}. \url{https://doi.org/10.4171/owr/2015/21}. The literature-source list labels this reference ``Oberwolfach Report 21.'' The supplied citation's year \(2016\) is retained despite the \(2015\) appearing in the DOI and proposed-year metadata.
\item Supplied literature reference: Shekhtman, \emph{On One Class of Hermite Projectors}. \url{https://doi.org/10.1007/s00365-015-9311-5}. No further bibliographic details are supplied.
\item Literature assessment date: \texttt{2026-08-21}. This is supplied dated curation, not independently verified. It labels the problem partially solved and reports a proof under an additional no-roots-at-infinity hypothesis while describing the unrestricted conjecture as unproved.
\item The registry labels the statement \texttt{corrected\_verified} and reports that the official Oberwolfach passage was checked directly and its definitions and notation restored. This is a supplied verification claim; no source visit or independent verification is asserted here.
\end{itemize}
\end{document}
